\section{Related Works}\label{app:related_work}

\textbf{Paradigms in CL.}
In recent years, interest in CL has risen, especially from the deep learning research community~\cite{achille2018life, nguyen2017variational, schwarz2018progress, tasnim2024recast}. CL approaches are generally categorized into three paradigms: \emph{regularization-based}, \emph{replay-based}, and \emph{architecture-based} methods \cite{parisi2019continual,nguyen2019toward,schwarz2018progress}.

Regularization-based methods constrain parameter updates to protect prior knowledge. \emph{Elastic Weight Consolidation} (EWC) \cite{kirkpatrick2017overcoming} and \emph{Synaptic Intelligence} (SI) \cite{zenke2017continual} penalize deviation from parameters important to previous tasks. \emph{Online EWC} \cite{chaudhry2018efficient} and \emph{Stable SGD}~\cite{mirzadeh2020dropout} extend these ideas to continuous task streams and smoother optimization.

Replay-based methods alleviate forgetting by explicitly revisiting prior data or learned representations. \emph{Experience Replay} (ER), \emph{Strong ER (SER)}~\cite{zhuo2023continual}, and \emph{Dark ER (DER)}~\cite{buzzega2020dark} maintain small exemplar buffers, while \emph{Generative Replay} techniques \cite{shin2017continual,van2018generative} synthesize pseudo-samples from previous distributions. Exemplar-based systems such as \emph{iCaRL} \cite{rebuffi2017icarl} combine replay with incremental classifiers.

Architecture-based approaches \cite{fernando2017pathnet,masse2018alleviating,li2017learning} mitigate interference by expanding or routing subnetworks to isolate knowledge across tasks. Although effective in preventing forgetting, they typically incur high computational or memory costs and reduced scalability. Despite their progress, these paradigms collectively continue to struggle with the \emph{stability–plasticity dilemma}, particularly under extended task sequences and non-stationary environments \cite{chen2023stability,de2021continual}.

\textbf{Multi-Timescale and Dual-Model Approaches.}
Beyond these standard paradigms, some recent works employ multiple learners operating at different temporal scales. DualNet~\cite{pham2021dualnet}, for example, maintains a fast learner and a slow learner inspired by Complementary Learning Systems. Their fast model is trained using replayed examples, combining both raw buffer samples and latent features from the slow model. The slow learner is also updated through replay. Although this architecture resembles our use of local and hierarchical models, the mechanisms differ substantially. Unlike dual-model systems which rely on standard rehearsal, HTCL repurposes the replay buffer to approximate the curvature of the loss landscape (Hessian) for a regularization step, rather than for blind rehearsal. Consequently, the hierarchical models evolve solely through a Taylor series approximation of the loss, enabling principled multi-timescale integration. This distinction allows HTCL to function as a model-agnostic consolidation layer rather than a replay-driven dual-network system.

\textbf{Task Ordering and Sequence Sensitivity.}
Recent studies demonstrate that task order profoundly affects CL performance \cite{bell2022effect,riemer2018learning,yoon2019scalable}. Models trained on identical tasks but different sequences can yield widely varying accuracies due to path-dependent optimization. While curriculum learning leverages structured ordering to encourage transfer, standard CL frameworks generally assume fixed or random task orders, lacking robustness under arbitrary permutations. Efforts to address this include additive parameter decomposition for order invariance \cite{yoon2019scalable} and order-robust objectives or loss regularizers \cite{dohare2024loss,guo2024continual}. Nonetheless, the factorial complexity of possible permutations makes brute-force optimization infeasible \cite{knoblauch2020optimal}, leaving order sensitivity an open challenge.

\textbf{Relation to Federated Learning.}
Our hierarchical design shares superficial similarity with multi-model frameworks in \emph{Federated Learning} (FL) \cite{mcmahan2017communication,li2020federated}. In FL, distributed clients train local models whose updates are aggregated on a central server to preserve privacy and reduce communication cost. FedAvg~\cite{mcmahan2017communication} aggregates local models by simple parameter averaging, implicitly assuming that all local updates are equally compatible. FedProx~\cite{li2020federated} extends this approach by introducing a proximal term that constrains local updates to remain close to the global model, improving stability under heterogeneous client distributions. FedCurv~\cite{shoham2019overcoming} employs curvature-based regularization to further stabilize aggregation.

However, CL differs fundamentally  from FL. FL addresses spatial heterogeneity across distributed clients, whereas CL handles temporal distributional shifts within a single learner. The curvature-aware consolidation in HTCL is conceptually related to methods like FedCurv, but HTCL integrates information hierarchically over time rather than across parallel clients. \textbf{Detailed experimental comparisons with FL baselines can be found in Appendix~\ref{app:federated}}, confirming that HTCL outperforms both FedAvg and FedProx when applied to sequential task integration and highlighting the importance of consolidation mechanisms tailored to temporal learning.

\textbf{Summary and Distinction.}
In summary, prior methods have mitigated forgetting via replay, regularization, or network expansion, but they remain vulnerable to task-order bias and scalability challenges. Few works explicitly formalize order robustness. In contrast, \emph{HTCL} introduces a principled, multi-level Taylor series framework that unifies local adaptability and hierarchical stability, yielding order-invariant and scalable CL behavior.